%!TEX root = main.tex

A \emph{sequence} over the set $X$ is $(a_0, \ldots, a_{n-1})$, where $a_i \in X$ for each $i \in [0,n-1]$, and $n$ is the \emph{length} of the sequence.
%A \emph{string} is a sequence over %$A$ is 
%a (finite) alphabet $\Sigma$. As usual, $\Sigma^*$ denotes the set of strings over $\Sigma$, and the empty sequence/string is denoted by $\varepsilon$.
We shall focus on \emph{string sequences}, namely, sequences over $\Sigma^*$. We use $()$ to denote the empty sequence.
% for some (finite) alphabet $\Sigma$. 

%We consider the following 
\smallskip
\noindent\emph{Operations on string sequences and strings.} %Let $\Sigma$ be an alphabet and 
We assume two string sequences $s_1 = (u_0, \ldots, u_{m-1})$ and  $s_2 = (v_0, \ldots, v_{n-1})$.
\begin{itemize}%[topsep=0pt,partopsep=0pt]
	\item $s_1 \concat s_2$ is the concatenation of $s_1$ and $s_2$, i.e., $(u_0, \ldots, u_{m-1}, v_0, \ldots, v_{n-1})$.

	      % \item $\transducer(u)$, which applies a \emph{functional} transduction $\transducer$ to $u\in \Sigma^*$ and returns  $v$ such that $(u, v) \in \transducer$. \tl{check}
	      %
	\item $\myat(s_1, i)$ is $u_i$ if $i\in [0,m-1]$, and  undefined otherwise.
	      %
	\item $\myjoin_v(s_1)$ joins the elements in $s_1$ with $v$ as the intermediate string, i.e., \polished{$\myjoin_v(s_1) = u_0 v u_1 v u_2 v \cdots v u_{m-1}$}. 
		      %
	\item $s_1[i \rightarrow u]$ replaces $u_i$ in $s_1$ by $u$ if $i\in [0,m-1]$,  and is $s_1$ otherwise,  i.e.
	\begin{center}
	 $[s_1[i \rightarrow u] = 
	\left \{
	\begin{array}{l l}
	(u_0, \ldots, u_{i-1}, u, u_{i+1}, \ldots, u_{m-1}), & \mbox{ if } i\in [0,m-1]; \\
	s_1, & \mbox{ otherwise}.
	\end{array}
	\right.
	$    
	\end{center}
	      %
	\item $\myfilter_e(s_1)$ filters out all elements in $s_1$ that do not match the regular expression $e$, i.e., $\myfilter_e(s_1) = (u_{i_1}, \ldots, u_{i_k})$  such that $0\le i_1 < \cdots < i_k\le m-1$ and  for all $j \in [0,m-1]$, $u_j \in \Lang(e)$ iff $j \in \{i_1, \ldots, i_k\}$. %$u_j \in \Lang(e)$ iff $j \in \{i_1, \ldots, i_k\}$. 
	      %
	      In particular, if none of $u_0, \ldots, u_{m-1}$ are in $\Lang(e)$, then $\myfilter_e(s_1) = ()$. 
	      %\tl{or $\myfilter_e(s_1) = ()$? it should be a sequence!?}\zhilin{yes, use () here}
%\tl{$\varepsilon$ or $(\varepsilon)$?}\zhilin{empty sequence, $\varepsilon$. $(\varepsilon)$ denotes a sequence of length 1, where the element is the empty string ?}
	      %
	\item $\mysplit_e(u)$ splits a string $u$ into a sequence of substrings according to the regular expression $e$, i.e.,  $\mysplit_e(u) = (u'_1, \ldots, u'_k)$, where $u = u'_1 v_1 u'_2 v_2 \cdots u'_{k-1}v_{k-1} u'_k$ such that (1) for each $i \in [k-1]$, $v_i$ is \polished{the leftmost and longest substring of $u'_i v_i u'_{i+1}\cdots v_{k-1}u'_k$} that matches $e$, and (2) $u'_k$ does not contain any substring that matches $e$.
	      % moreover, $u'_k$ does not contain occurrences of $v$.
	      In particular, if $u$ does not contain any substring that matches $e$, then $\mysplit_e(u) = (u)$.

	\item $s_1[i, j]$ is defined as follows: if $i\in [0,m-1]$ and $j \ge 1$, then $s_1[i, j]$ is  the subsequence $(u_i, \ldots, u_{\min(i+j-1,m-1)})$; if $i\in [0,m-1]$ and $j = 0$, then $s_1[i, j]$ is the empty sequence $()$;  if $i\not\in [0,m-1]$ or $j < 0$, then $s_1[i, j]$ is undefined. 
	%\tl{can we try to avoid $\bot$? tbd}

	      % \item $\mymatch_e(u)$ extracts the leftmost occurrence of $e$ in a string $u$ and returns the sequence $(v)$, where $u = u'_1 v u'_2$ such that $v$ is the leftmost occurrence of $e$ in $u$.  
	      %
	\item $\mymatchall_e(u)$ extracts the substrings of $u$ that match the regular expression $e$, 
	%all occurrences of $e$ in the string $u$,
	 i.e., $\mymatchall_e(u) = (v_1, \ldots, v_{k-1})$ such that $u = u'_1 v_1 u'_2 v_2 \cdots u'_{k-1}v_{k-1} u'_k$ and (1) for each $i \in [k-1]$, $v_i$ is \polished{the leftmost and longest substring of $u'_i v_i u'_{i+1}\cdots v_{k-1}u'_k$}, and (2) $u'_k$ has no substring that matches $e$.

	      % \item $\transducer(s_1)$ applies a transduction $\transducer$ to $s_1$ elementwise, i.e., $(\transducer(u_1), \ldots, \transducer(u_m))$.
	      %subsequence, match, matchAll, transduction, 
\end{itemize}

\polished{Note that we choose the longest match semantics in $\mysplit_e(u)$ and $\mymatchall_e(u)$
just for illustrating our approach, where
the shortest match semantics can be captured by adapting the automata construction in the pre-image computation.}

%A string over an alphabet $\Sigma$ is a sequence of symbols from $\Sigma$. 
\smallskip
\noindent\emph{The logic $\arraylogic$ of string sequences.} $\arraylogic$   
has three data types: $\inttype$ (integer), $\strtype$ (strings), and $\seqtype$ (sequences).
%In this paper, $\seqtype$ is seen as a composite data type \tl{not sure about this term} where the elements are of type $\strtype$. 
%
We use $u, v, \ldots$ (resp. $\svaru, \svarv, \ldots$) to denote string constants (resp. variables); $m,n, \ldots$ (resp. $\svarm, \svarn, \ldots$) to denote integer constants (resp. variables);
$s, t, \ldots$ (resp. $\svars,\svart, \ldots$) to denote sequence constants (resp. variables); and $e$ to denote regular expressions.
%We define a theory $\arraylogic$ with three sorts: \textbf{string}, \textbf{integer} and \textbf{array}. The \textit{elements} of \textbf{array} are in \textbf{string} sort, and the $index$ of \textbf{array} is in \textbf{Integer} sort. 
%\begin{figure}[htbp]
%	\renewcommand{\arraystretch}{1.3}

The syntax of $\arraylogic$ is defined as follows:
\[
\begin{tabular}{c l l l r}
		%        $u, v, \cdots$     &  &                                                                          & string constants           \\
		%        $\svarx, \svary, \cdots$     &  &                                                                          & string variables           \\
		%        $m,n, \cdots$     &  &                                                                          & integer constants         \\
		%        $\ivarx, \ivary, \cdots$     &  &                                                                          & integer variables          \\
		%        $s, s', \cdots$     &  &                                                                          & sequence constants            \\
		%        $\seqvara,\seqvarb, \cdots$     & &                                                                          & sequence variables            \\
	$~$ &	$it$      & $ \eqdef $ & $m \mid \svarm \mid it+it \mid it-it \mid \mystrlen(strt)\mid \myseqlen(seqt) $                                                      & integer terms  \\
		&	$strt$    & $ \eqdef $ & $u \mid \svaru \mid strt \concat strt \mid \myat(seqt, it) \mid \myjoin_u(seqt)    $                      & string terms   \\
		&	$seqt$    & $ \eqdef $ & $(strt) \mid s \mid \svars \mid seqt \concat seqt \mid seqt[it \rightarrow strt] \mid  \myfilter_e(seqt) \mid  $ &                \\
			&          &            & $\ \ seqt[it, it]  \mid \mysplit_e(strt) \mid \mymatchall_e(strt) $                              & sequence terms \\
		%        $a$      & $::=$    & $\seqvara = \seqvarb[m: n] \mid A = \transducer(B) \mid x=\myread(A, m) \mid$          &                           \\
		%                  &          & $A = \mywrite(B, x ,m) \mid x = \myjoin(A,u) \mid A = \mysplit(x, e) \mid$ &                           \\
		%                  &          & $x= y\cdot z \mid x = \transducer(y)\mid x\in e$                         & atomic formulas        \\
		&	$\varphi$ & $ \eqdef $ & $ it \iop it  \mid seqt = seqt \mid strt = strt \mid strt \in e  \mid \varphi \wedge \varphi $                   & formulas
	\end{tabular}
    \]
where  $\iop\, \in \{=, \neq, <, \le, >, \ge\}$.
%	\caption{Syntax of $\arraylogic$}
%	\label{fig:syntax}
%\end{figure}
%
%The syntax of $\arraylogic$ is defined by the BNF in Figure~\ref{fig:syntax}, 

Here, $it$, $seqt$ and $strt$ denote integer, sequence and string terms, respectively. In particular,
the integer terms are linear arithmetic expressions where $\myseqlen(\svars)$ (resp. $\mystrlen(\svaru)$) denotes the length of a sequence $\svars$ (resp. string $\svaru$);
the sequence terms are constructed from a singleton sequence, %(sequences of length one), 
sequence constants and sequence variables by applying the aforementioned sequence/string operations.
%: concatenation ($\concat$), write ($seqt[it \rightarrow strt]$, writing a string $strt$ to a position $it$ in the sequence $seqt$), split operation ($\mysplit(strt, u)$),  subsequence ($seqt[it, it]$), $\mymatch(strt, e)$, $\mymatchall(strt, e) $, and $\transducer(seqt)$. 
%
A $\arraylogic$ formula is a conjunction of atomic formulas, each of which is of the form $it_1 \iop it_2$ where $\iop\, \in \{=, \neq, <, \le, >, \ge\}$,  a sequence equality $seqt_1 = seqt_2$, or a string equality $strt_1 = strt_2$.
%
%For each regular expression $e$, the language defined by $e$, denoted by $\Lang(e)$, is defined recursively. For instance, $\Lang(\emptyset) = \emptyset$, $\Lang(\varepsilon) = \{\varepsilon\}$, $\Lang(a) = \{a\}$, and $\Lang(e_1 \cdot e_2) = \Lang(e_1) \cdot \Lang(e_2)$, $\Lang(e_1 + e_2) = \Lang(e_1) \cup \Lang(e_2)$, and so on. 
%
Note that sequence and string inequalities can be expressed as well. %seen as macros, with the help of existential quantifiers. 
For instance (using disjunctions and quantifiers for the sake of presentation):
%
%$$\svarx \neq \svary ~\equiv~ \exists \svarx_1, \svarx_2, \svary_2.\ \bigvee \limits_{a, b \in \Sigma, a \neq b} (\svarx = \svarx_1 a \svarx_2 \wedge \svary = \svarx_1 b \svary_2) \vee \exists \svarz.\ \bigvee_{a \in \Sigma} (\svarx = \svary a \svarz \vee \svary = \svarx a \svarz),$$
%\tl{how about $\varepsilon$ and $a$?}
%\tl{does not read right ?}%
 \begin{align*}
    \svarx \neq \svary & \equiv& \exists \svarx_1, \svarx_2, \svary_2. \bigvee \limits_{a, b \in \Sigma, a \neq b} (\svarx = \svarx_1 a \svarx_2 \wedge \svary = \svarx_1 b \svary_2) \vee \exists \svarz. \bigvee_{a \in \Sigma} (\svarx = \svary a \svarz \vee \svary = \svarx a \svarz) \\
		\seqvara \neq \seqvarb &\equiv & (\exists \seqvara_1, \seqvara_2, \seqvarb_2, \svarz_1, \svarz_2.\ \seqvara = \seqvara_1 \concat (\svarz_1) \concat \seqvara_2 \wedge \seqvarb = \seqvara_1 \concat (\svarz_2) \concat \seqvarb_2 \wedge \svarz_1 \neq \svarz_2)\ \vee \\
		                   &   & %\hspace{4mm} 
                               \exists \seqvarg, \svarz.\ \seqvara  = \seqvarb \concat (\svarz) \concat \seqvarg \vee    \seqvarb  = \seqvara \concat (\svarz) \concat \seqvarg.
\end{\align*}
%
\polished{Though rewriting inequalities produces disjunctions, they can still be handled in the DPLL(T) framework~\cite {dpll_t}. Hence, we focus on the disjunction-free fragment.}

%%%%%%%%%%%%%%%%%%%%%%%%%%%%%%
%\hide{
%	A $\arraylogic$ formula $\varphi$ is a conjunction of formulas of the form $t_1\circ t_2$, $A = B[m:n]$, $A = \transducer(B)$, $x = \myread(A,m)$, $A = \mywrite(B, x, m)$, $x=\myjoin(A,u)$, $A=\mysplit(x,e)$, $x = \transducer(y)$, $x=y\cdot z$ or $x\in e$. Definitely, $t$ is a integer term built from the constant $m$, the varible $i$, the size of array $\mylen(A)$ and the length of string $\mylen(x)$ by operators $+$ and $-$; $\circ \in \{=,\not=, <, \leq, >, \geq \}$ is the set of arithmatic binary predicates on integer terms $t_1$ and $t_2$; $e$ is a regular expression; $B[m:n]$ is the sub-array of $B$ containing elements between $m$ and $n$; $\transducer(B)$ modifies the elements of $B$ by transducer $\transducer$ and returns the result; $\myread(A, m)$ returns the $m$-th element of $A$; $\mywrite(B, x ,m)$ returns the array $B$ after $x$ is written to $m$-th position of $B$; $\myjoin(A,u)$ joins all of the elements in $A$ by string $u$; $\mysplit(x,e)$ splits $x$ into an ordered string array by regex $e$; $\transducer(y)$ returns the output of the transducer $\transducer$ running on $y$; $y\cdot z$ is the concatenation of $y$ and $z$.
%}
%%%%%%%%%%%%%%%%%%%%%%%%%%%%%

\begin{proposition}\label{prop-undec}
	The satisfiability of $\arraylogic$ is undecidable.
\end{proposition}
%Idea: replaceAll where the replacement strings are constants can be captured by transducers. 
%
%From Proposition~\ref{prop-undec}, we know that the satisfiability of $\arraylogic$ is undecidable. In the sequel, we show that for the straight-line fragments of $\arraylogic$, the satisfiability is decidable. 

We define the \emph{straight-line fragment} $\slarraylogic$ of $\arraylogic$. It is easy to see that by introducing fresh %sequence/string 
variables, we can rewrite a $\arraylogic$ formula to a form in which the left-hand side of each equality $seqt_1 = seqt_2$ (resp. $strt_1 = strt_2$) is a sequence (resp. string) variable.
%
%i.e., the string/sequence equalities are in one of the following forms,
%\[
%	\begin{array} {l}
%		\svars = (\svaru) \mid \svars = s \mid \svars = \svart \mid \svars = \svart_1 \concat \svart_2 \mid \svars = \svart[it \rightarrow \svaru] \mid \\
%		  \svars  = \myfilter_e(\svart) \mid \svars = \svart[it_1, it_2] \mid \svars = \mysplit_e(\svaru) \mid \svars = \mymatchall_e(\svaru) \mid \\
%		\svaru = u \mid \svaru = \svarv \mid \svaru = \svarv_1 \concat \svarv_2 \mid \svaru = \myat(\svars, it) \mid \svaru = \myjoin_u(\svars).
%	\end{array}
%\]
Note that the original and rewritten formulas are equisatisfiable. %rewriting preserves the satisfiability. 
Henceforth, we assume that all sequence/string equalities satisfy this constraint unless otherwise stated. %of the aforementioned forms. 

%It is easy to see that for a sequence equality $seqt_1 = seqt_2$, we can introduce a fresh sequence variable $\seqvara$ and rewrite  it as $\seqvara= seqt_1 \wedge \seqvara = seqt_2$. Similarly, for a string equality $strt_1 = strt_2$, we can introduce a fresh string variable $\svarx = strt_1 \wedge \svarx = strt_2$. Note that the rewritings preserve the satisfiability. As a result, we can assume that all the sequence equalities $seqt_1 = seqt_2$ (resp. string equalities $strt_1 = strt_2$) satisfy that $seqt_1$ is a sequence variable (resp. $strt_1$ is a string variable). 

\begin{definition}[Straight-line fragment]\label{def:straight_line}
	Given a $\arraylogic$ formula $\varphi$, 
	 the \emph{dependency graph} $G_\varphi$ of  $\varphi$
	 is a directed graph $(V,E)$ such that
	%	we define the \emph{dependency graph} $G_\varphi$ of  $\varphi$, such that 
	$V$ is  the set of string and sequence variables in $\varphi$, and % the graph comprising 
	 $(v_1, v_2)\in E$ %such that $v_1,v_2$ are sequence or string variables, 
	iff  $v_1 = {\tt rhs}$ is a sequence or string equality in $\varphi$ and $v_2$ occurs in ${\tt rhs}$ except for integer terms.
	%
	%For instance, from the equality $\seqvara = \seqvarb[1, \mylen(\svarx)]$, we know that $(\seqvara, \seqvarb)$ is an edge in $G_\varphi$, while $(\seqvara, \svarx)$ is not. 

$\varphi$ is \emph{straight-line} if $\varphi$ is disjunction-free, and each string/sequence variable occurs as the left-hand side of equalities at most once; moreover, $G_\varphi$ is acyclic.

%	$\varphi$ is \emph{straight-line} if $G_\varphi$ is acyclic.
\end{definition}
%Let us use $\slarraylogic$ to denote the straight-line fragment of $\arraylogic$.

\begin{example}
$\svars = \mysplit_e(\svaru) \wedge \svart = \svars[1, \myseqlen(\svars) - 1]$ is  straight-line, while  $\svars = \svars[1, \myseqlen(\svars) - 1]$ is not, since there is a self-dependency on $\svars$.
\end{example}

%Straight-line fragments: Write the equalities as $\svarx = strt$ or $\seqvara = seqt$. The dependency graph should be acyclic (where the string variables and sequence variables are all taken into consideration). 

%\zhilin{stopped here}

%%%%%%%%%%%%%%%%%%%%%%%%%%%%%%%%%%
%%%%%%%%%%%%%%%%%%%%%%%%%%%%%%%%%%
\hide{
	%\subsection{Definition of Straight-Line Formula }
	\begin{definition}[Relational constraints ]
		Relational constraints $\phi$ are defined by the following rules:
		\begin{align*}
			\phi \ ::= \  & A = B[m:n] \mid A = \transducer(B) \mid x=\myjoin(A,u) \mid A = \mysplit(x,e) \mid                       \\
			              & x=\myread(A, m) \mid A = \mywrite(B, x ,m) \mid x = \transducer(y) \mid x=y\cdot z \mid \phi \wedge \phi
		\end{align*}
	\end{definition}
	\begin{definition}[Straight-line relational constraints]
		A relational constraint $\phi$ is straight-line if $\phi = \bigwedge\limits_{1\leq i \leq m} \chi_i = P_i$ such that
		\begin{itemize}
			\item $\chi_1,\cdots,\chi_m $ are mutally distinct string varibles and array varibles.
			\item For each $i\in [m] $, all the varibles in $P_i$ are in the set $\{\chi_1,\cdots,\chi_{m-1} \}$.
		\end{itemize}
	\end{definition}
	\begin{definition}[Straight-line formula]
		A formula $F$ of $\arraylogic$ is straight-line iff all relational constraints in $F$ are straight-line.
	\end{definition}
}
%%%%%%%%%%%%%%%%%%%%%%%%%%%%%%%%%%
%%%%%%%%%%%%%%%%%%%%%%%%%%%%%%%%%%